\subsection{Even centers of minimum two and unit-arm reductions}\label{loc:S38}
In this proof $D_j=\delta_{j}/(1-\theta^2)$ bounds finite same-sign tails.
For $2=a_1<b_1<a_2<\cdots<a_t<b_t$ use
$F=f_2h_{b_1}\prod_{i=2}^th_{a_i}f_{b_i}$.
We prove $\E[q^3F^4]\le e^{-d/20}$, with the stronger loss
$[8(t-2)+1]d$ when $t\ge2$.
In particular the first pair is $f_2h_b$, rather than the terminal
direction $h_2f_b$.

\proofpart{A logarithmic comparison retaining both arms}
Suppose $h_b\le\widehat h\le2$ and
$f_2+\widehat h=2-\Delta+R$, $R=KV$.
Write $I=\|V\|^2$, $J=\|R\|^2$,
$t_0=\E[gV]$, $L_0=\E[rV]$. Then
\begin{equation}\label{eq:compact-min2-center-master}
 \E[q^3(f_2h_b)^4]\le e^\Xi,
\end{equation}
where the complete exponent is
\[
 \Xi=-4\Delta+8091d/2056+900I/257+6J+9L_0
                         +(-308/257+2\Delta)t_0+\Delta^2/6.
\]
To prove it put $u=f_2-1=g+r$, $v=\widehat h-1$,
$w=(u-v)/2=g+r-R/2+\Delta/2$.
The quadratic logarithmic bound gives
\[
 q^3(f_2\widehat h)^4\le
 e^{-4\Delta-3g+4R-9g^2/10-12w^2/5}.
\]
Take $X=-3g+4R-2\Delta/3=-5g/3-4w/3+KZ$,
$Z=(4g+10V)/3$.
Weighted Cauchy and the original row bound give
$X^2/7\le9g^2/10+12w^2/5+\eta$,
$\eta=(9/257)\|4g+10V\|^2$.
The scalar inequality $e^{X-X^2/7}\le Q(X)=1+X+3X^2/8$
follows from derivative $X(6X^2-5X+2)/(56Q)$.
Finally
\[
 \eta=144d/257+900I/257+720t_0/257,\quad
 \E X=-4t_0-2\Delta/3,
\]
and $\E X^2=9d+16J+24L_0+(16\Delta/3)t_0+4\Delta^2/9$.
Taking expectation under the fixed probability measure gives the exponent.

\proofpart{Every single pair}
For $b=3$ take $\widehat h=h_3$, $\Delta=d$, $R=-s$, $V=-r$.
Thus $(I,J,t_0,L_0)=(\delta_{2},\delta_{3},0,-\delta_{2})$ and
\[
 \Xi=-133d/2056-1413\delta_{2}/257+6\delta_{3}+d^2/6\le-d/16.
\]
For $b\ge4$, use the retained coefficients $C_1,C_2,C_3$ from
Section~\ref{loc:S28} and
$\widehat h=2-\beta_b-S_b-C_1h_1-C_2h_2-C_3h_3$.
It lies between the actual $h_b$ and two.
Let $R_0=KV_0$ be the same first-pair polynomial as there.
Then $\Delta=2d$, $R=R_0-s$, $V=V_0-r$.
Put $M=\E[gV_0]$, $N_0=\E[gR_0]$, $N_1=-\E[sR_0]$.
The quantities entering the exponent are exactly
\[
 t_0=M,\quad L_0=-N_0-\delta_{2},\quad
 I=\|V_0\|^2+\delta_{2}+2N_0,\quad J=\|R_0\|^2+\delta_{3}+2N_1.
\]
Hence
\begin{align*}
 \Xi={}&-8357d/2056+(900/257)\|V_0\|^2+6\|R_0\|^2
 -1413\delta_{2}/257-513N_0/257\\
 &+6\delta_{3}+12N_1+(-308/257+4d)M+2d^2/3.
\end{align*}
For $b=4$, the exact inputs are
\[
 I=3d^2+\delta_{2}+\delta_{3},\quad J=d^3+\delta_{3}+\delta_{4}+2d\delta_{2},
 \quad t_0=\delta_{2},\quad L_0=-d^2-\delta_{2}.
\]
The linear spectral terms are at most $-5\delta_{2}/3$;
using $d\le1/3$, $\delta_{2}\le5d/12$ gives $\Xi\le-19d/18$.
For $b=5$, they are
\begin{align*}
 I&=3d^2+\delta_{2}+2d\delta_{2}+(1+d)\delta_{2}^2-\delta_{3}+\delta_{4},\\
 J&=d^3+2d^2\delta_{2}+d\delta_{2}^2+\delta_{2}^3+(1-2d)\delta_{3}-\delta_{4}+\delta_{5}
                                      +2d\delta_{2}+2\delta_{2}^2,\\
 t_0&=(1+d)\delta_{2},\qquad L_0=-d^2-d\delta_{2}+\delta_{3}-\delta_{2}.
\end{align*}
Substitution bounds the exponent by
\[
 -4d+9d^2/4+6d^3-2\delta_{2}+8d\delta_{2}+16d^2\delta_{2}+(16+10d)\delta_{2}^2+6\delta_{2}^3.
\]
The positive costs divided by $d$ are at most
$3/4+2/3+55/54+29/27+7/150$, giving
$\Xi\le-1199d/2700\le-2d/5$.

For every $b\ge6$, the common finite-frequency formulas give
\begin{gather*}
 \|V_0\|^2\le C^2+(1+\theta)\delta_{3}+D^2d,\quad
 \|R_0\|^2\le C^2d+\delta_{4}+\delta_{5}+D^2\delta_{2},\\
 N_0\ge Cd-\delta_{3}-E\delta_{2},\quad
 N_1\le C\delta_{2}-\delta_{4}+\delta_{5}-E\delta_{3}.
\end{gather*}
For $d\le1/5$, the $M$ coefficient is negative and $M\ge0$.
Let $A=900/257$, $B=513/257$.
The higher spectral terms after substitution have nonpositive sum:
their coefficient per $\delta_{2}$ is at most
$(41395/3200)(13/32)+7/32-549/100=-1641/102400$.
The remaining positive terms are bounded as follows, using
$d\le C\le257d/160$, $D\le81d/80$, $E\le D$:
\[
\begin{array}{c|c}
\text{term divided by }d&\text{upper bound}\\\hline
[(A+6d)C^2-BCd]/d&9/5\\
D^2(Ad+6\delta_{2})/d&1/4\\
12C\delta_{2}/d&8/5\\
BE\delta_{2}/d&1/6\\
2d/3&2/15
\end{array}
\]
Their sum is $79/20$, while $8357/2056>4$.
Thus $\Xi\le-d/20$ on the full small interval, with every finite $b$.
For $d\ge1/5$, the stronger actual path means
\eqref{eq:compact-strong-path-means} give $h_b\le77/80$.
The same-direction singleton result of Section~\ref{loc:S36} gives
$\E[q^3f_2^4]\le1$, so the actual pair fourth moment is at most
$(77/80)^4\le e^{-3/20}\le e^{-d/20}$.

\proofpart{All later pairs in the same integral}
Let $S=KW=\sum_{i=2}^t\sum_{j=a_i}^{b_i-1}K^jg$ and put
$I_W=\|W\|^2$, $J_W=\|S\|^2$,
$t_W=\E[gW]$, $L_W=\E[rW]$.
In the preceding logarithmic argument replace $X$ by $X+4S$ and
the row vector $Z$ by $Z+4W$.
The exact additional cost is
\[
 \mathrm{Extra}=2160\langle V,W\rangle/257+1296I_W/257
 +12\langle R,S\rangle+6J_W+9L_W+(-164/257+2\Delta)t_W.
\]
The total exponent is $\Xi_{\rm first}-8(t-1)d+\mathrm{Extra}$.
On $d\le1/5$, $|-164/257+2\Delta|<2/3$,
$|t_W|\le D_3<d/5$ and $L_W\le D_4<21d/250$.
For $b=3$ use
$\|V\|\le13\sqrt d/20$, $\|R\|\le5\sqrt d/12$,
$\|W\|<2\sqrt d/5$, $\|S\|<\sqrt d/4$.
The six term cost ratios sum to
$221/100+102/125+5/4+3/8+189/250+2/15=16621/3000<7$.
For $b\ge4$, the exact norms above imply
$\|V\|\le7\sqrt d/5$, $\|R\|\le5\sqrt d/6$.
Indeed $C\le13d/8$, $D\le21d/20$ give squared ratios below
$23/12<49/25$ and $1333/2000<25/36$.
Later minimum at least five gives
$\|W\|<\sqrt d/4$, $\|S\|<13\sqrt d/80$;
the six costs are at most
$119/40+51/160+13/8+507/3200+189/250+2/15<7$.
This proves the stronger multiple-pair loss $[8(t-2)+1]d$.
For $d\ge1/5$, the actual tail is bounded by $S<2\sqrt d/5\le d$.
Multiply the same-direction first-pair bound by its actual exponential
tail, giving the still stronger loss $[8(t-2)+4+1/20]d$.

\proofpart{A residual unit arm}
Except for the sole residual pair $(1,2)$, direct that unit arm outward.
If $G$ is the remaining actual message with $\E[pG^4]\le1$, then
\[
 \E[q^3(pG)^4]=\E[p(pq)^3G^4]\le\E[pG^4]\le1.
\]
For sole $(1,b)$, $b\ge3$, take $G=\bar f_{b}$ and use
Section~\ref{loc:S20}. If at least three arms remain after removing the
unit arm, use the fully reversed positive odd canonical direction
from Sections~\ref{loc:S15}--\ref{loc:S16}.
These are the original zero-step fourth directions.
The exception $(1,2)$ has length outside the path theorem and is handled
separately in the main star catalogue.
All equal pairs, zero parameters and full reversal follow from the common rules.
